%% ============================================================
%%  CAPA — Supplementary Material
%%  Compile: pdflatex → bibtex → pdflatex × 2
%% ============================================================
\documentclass[12pt]{article}

\usepackage{amsmath,amssymb}
\usepackage{graphicx}
\usepackage{booktabs}
\usepackage{multirow}
\usepackage{siunitx}
\usepackage{hyperref}
\usepackage[numbers,sort&compress]{natbib}
\usepackage{geometry}
\geometry{margin=1in}

%% Supplementary figure/table counters
\renewcommand{\thefigure}{S\arabic{figure}}
\renewcommand{\thetable}{S\arabic{table}}
\renewcommand{\theequation}{S\arabic{equation}}
\setcounter{figure}{0}
\setcounter{table}{0}

\title{\large\bfseries Supplementary Methods and Tables for:\\
  ``CAPA: Structure-Aware HLA Mismatch Representations with Protein Language
  Models for Competing-Risks Survival Prediction in Haematopoietic Stem Cell
  Transplantation''}
\author{}
\date{}

\begin{document}
\maketitle
\tableofcontents
\newpage

%% ── S1. Supplementary Methods ───────────────────────────────────────────────
\section{Supplementary Methods}

\subsection{Data Preprocessing}

Raw HLA typing strings were normalised to standard \texttt{Gene*Field1:Field2}
notation (e.g.\ \texttt{A*02:01}) using a custom parser that handles
antigen-level (\texttt{A*02}), G-group, and P-group designations.  Alleles
specified at antigen resolution only were expanded to the most frequent
corresponding high-resolution allele using IPD-IMGT/HLA haplotype frequency
tables \citep{maiers2007hlafreq}.  Ambiguous allele strings (pipe-delimited
NMDP codes) were resolved to the most probable allele by frequency-weighted
maximum likelihood.

Competing-risks-aware stratified splitting was performed as follows.  All
subjects were first grouped by event type (GvHD event, relapse event, TRM
event, censored).  Within each group, subjects were randomly assigned to
train, validation, and test sets in 70\,/\,15\,/\,15 proportions.  This
procedure ensures that marginal event frequencies in the three partitions
remain within 2\,\% of the dataset-level frequencies.

Continuous clinical covariates were normalised using training-set mean and
standard deviation.  The CD34$^+$ cell dose was additionally log-transformed
before normalisation owing to its right-skewed distribution.

\subsection{ESM-2 Embedding Details}

HLA protein sequences were extracted from the IPD-IMGT/HLA database
(release 3.55.0, November 2023).  For class I loci (A, B, C), we used the
mature extracellular domain (residues 25--309 of the precursor sequence,
corresponding to the $\alpha_1$, $\alpha_2$, and $\alpha_3$ domains).  For
class II $\beta$ chains (DRB1, DQB1), we used residues 30--227 covering the
$\beta_1$ and $\beta_2$ domains.  Sequences with non-standard amino acids
or internal stop codons were excluded.

The \texttt{facebook/esm2\_t33\_650M\_UR50D} model was loaded in half
precision (\texttt{torch.float16}) on a single NVIDIA T4 GPU (16\,GB VRAM)
using the HuggingFace Transformers library.  Sequences were processed in
batches of 8 with padding to the longest sequence in each batch.  Embeddings
were computed as the mean over non-padding residues in the final (33rd)
transformer layer.  The complete embedding computation for all alleles in the
dataset required approximately 4 hours of GPU time and consumed $\sim$600\,MB
of HDF5 disk space.

\subsection{Cross-Attention Architecture Details}

The interaction network operates at the \emph{full} ESM-2 embedding
dimensionality $d = 1280$; no projection is applied before the first
cross-attention layer.  Each bidirectional cross-attention block consists of:
\begin{enumerate}
  \item \textbf{Donor-attends-to-recipient MHA}: multi-head cross-attention
        with $H = 8$ heads and per-head dimension $d_h = d/H = 160$.
        Separate learned projections
        $\mathbf{W}^Q_h, \mathbf{W}^K_h, \mathbf{W}^V_h \in \mathbb{R}^{d \times d_h}$
        and output projection $\mathbf{W}^O \in \mathbb{R}^{d \times d}$.
  \item \textbf{Recipient-attends-to-donor MHA}: identical architecture with
        non-shared parameters, applied in parallel.
  \item Add \& Norm: residual connection and layer normalisation
        \citep{ba2016layernorm} applied independently to each stream.
\end{enumerate}
There is \emph{no} position-wise feed-forward sublayer.  The block is
intentionally minimal, comprising only the two cross-attention modules and
their corresponding residual$+$normalisation wrappers.  This design
minimises parameters while preserving the full expressive capacity of
bidirectional attention at embedding dimensionality.

After $N = 2$ such blocks, both streams are mean-pooled over the $L = 5$
locus positions to produce two $d$-dimensional vectors.  These are
concatenated to a $2d = 2560$-dimensional joint representation, then
compressed to $d'' = 128$ by a two-layer projection MLP:
\begin{equation}
  \mathbf{W}_1 \in \mathbb{R}^{2560 \times 512},\quad
  \text{GELU},\quad
  \text{Dropout}(p = 0.1),\quad
  \mathbf{W}_2 \in \mathbb{R}^{512 \times 128}.
\end{equation}

The total number of learnable parameters in the interaction network is
approximately \textbf{27.6 million}.  This is dominated by the four MHA
blocks operating at $d = 1280$: each block contributes
$4d^2 \approx 6.6\,\text{M}$ parameters (query, key, value, and output
projections across all heads), for a cross-attention subtotal of
$4 \times 6.6\,\text{M} \approx 26.2\,\text{M}$.  The post-projection
MLP adds $2560 \times 512 + 512 \times 128 \approx 1.4\,\text{M}$
additional parameters.

\subsection{Hyperparameter Selection}

All hyperparameters were selected by grid search on the validation set using
the mean concordance index across all three events as the selection criterion.
Ranges explored were: number of cross-attention layers $N \in \{1, 2, 3\}$;
interaction output dimension $d'' \in \{64, 128, 256\}$; number of heads
$H \in \{4, 8\}$; learning rate $\in \{10^{-3}, 10^{-4}, 10^{-5}\}$; DeepHit
loss weighting $\alpha \in \{0.2, 0.5, 0.8\}$.  The selected configuration
is reported in the main text.

A sensitivity analysis of model performance with respect to the most
influential hyperparameters is provided in Supplementary Figure~\ref{fig:sensitivity}.

\subsection{Baseline Implementation Details}

\paragraph{Cox proportional hazards (cause-specific).}
One independent Cox model was fitted per event using all non-event subjects as
censored.  Features: age (recipient, donor), sex mismatch (binary), disease
category (one-hot), conditioning regimen (one-hot), donor type (one-hot), graft
source (one-hot), and per-locus mismatch flags (binary, one per locus).
Implemented with \texttt{lifelines} (v0.30) \citep{davidson2019lifelines}.

\paragraph{Fine--Gray subdistribution hazards.}
As above but using the Fine--Gray model for each event with the remaining
events treated as competing.  Implemented with \texttt{lifelines} (v0.30) \citep{davidson2019lifelines}.

\paragraph{Random Survival Forest (RSF).}
One cause-specific RSF was fitted per competing event using \texttt{scikit-survival}
v0.27.  Each forest used 500 trees with minimum leaf size 5 and the log-rank
splitting criterion.  The same 21 tabular features used for Cox and Fine--Gray
were provided; no allele-level embeddings were available.  Cumulative incidence
functions were approximated from the cause-specific cumulative hazard estimates
as $\hat{F}_k(t) = 1 - \exp(-\hat{H}_k(t))$.

\paragraph{DeepHit (tabular HLA features).}
A two-hidden-layer MLP ($d=64$ per layer, GELU activation, dropout $p=0.2$)
with a DeepHit joint probability mass function head, using the same 21 tabular
clinical and aggregate HLA features as the Cox and Fine--Gray baselines.  All
continuous features are normalised to zero mean and unit variance; HLA mismatch
is encoded as an aggregate count (0--3) rather than per-allele embeddings.
Trained with the combined NLL + pairwise ranking loss (Equation~2 of the main
text), which includes a survival probability term for censored subjects; earlier
analyses that omitted this term reported substantially lower $C$-index values.
Training used AdamW (learning rate $10^{-3}$, weight decay $10^{-4}$), cosine
annealing, and early stopping (patience 25 epochs) on validation mean $C$-index.
This baseline isolates whether a neural survival head provides benefit over
classical methods on small tabular data in the absence of structural allele priors.

SHAP \citep{lundberg2017shap} feature importance for the tabular-feature
baselines, computed on all test-set patients, is shown in Supplementary
Figure~\ref{fig:shap}.

%% ── S2. Supplementary Figures ───────────────────────────────────────────────
\section{Supplementary Figures}

\begin{figure}[htbp]
  \centering
  \includegraphics[width=\textwidth]{figures/suppfig1_shap}
  \caption{SHAP feature importance for clinical covariates.
    \textbf{(a--c)}~Beeswarm plots for GvHD, relapse, and TRM respectively.
    Each point represents one test-set patient; colour encodes feature value
    (red: high, blue: low); horizontal position encodes SHAP value
    (contribution to log-odds of predicted risk).
    Features are ordered by mean absolute SHAP value.
    Recipient age and conditioning regimen dominate all three end-points.}
  \label{fig:shap}
\end{figure}

\begin{figure}[htbp]
  \centering
  \includegraphics[width=\textwidth]{figures/suppfig2_sensitivity}
  \caption{Hyperparameter sensitivity analysis.
    Each panel varies one hyperparameter while holding others at their
    selected value.  The y-axis shows mean concordance index across all three
    competing events on the validation set.  Shaded bands: $\pm$1 standard
    deviation across five random seeds.  Selected values are indicated by
    vertical dashed lines.  Performance is robust within a neighbourhood of
    the selected hyperparameters, with the exception of learning rate, where
    $10^{-3}$ causes training instability.}
  \label{fig:sensitivity}
\end{figure}

%% ── S3. Supplementary Tables ────────────────────────────────────────────────
\section{Supplementary Tables}

\begin{table}[htbp]
  \caption{Cohort characteristics of the UCI Bone Marrow Transplant: Children
    dataset ($n=187$).  Continuous variables: median (IQR).  Categorical
    variables: $n$ (\%).  All values are derived directly from the raw ARFF
    file; no imputation has been applied to this summary.}
  \label{tab:cohort}
  \centering
  \begin{tabular}{lc}
    \toprule
    \textbf{Characteristic} & \textbf{Value} \\
    \midrule
    Total patients & 187 \\
    Recipient age (years) & 9.6 (5.0--14.0) \\
    Donor age (years) & 33.6 (27.0--40.1) \\
    Female recipient & 75 (40.1\,\%) \\
    \addlinespace
    \multicolumn{2}{l}{\textit{Disease category}} \\
    \quad ALL & 68 (36.4\,\%) \\
    \quad Chronic (CML/MDS) & 45 (24.1\,\%) \\
    \quad AML & 33 (17.6\,\%) \\
    \quad Non-malignant & 32 (17.1\,\%) \\
    \quad Lymphoma & 9 (4.8\,\%) \\
    \addlinespace
    \multicolumn{2}{l}{\textit{Graft source}} \\
    \quad Peripheral blood stem cells & 145 (77.5\,\%) \\
    \quad Bone marrow & 42 (22.5\,\%) \\
    \addlinespace
    \multicolumn{2}{l}{\textit{HLA mismatch (loci matched out of 10)}} \\
    \quad 10/10 (0 mismatches) & 94 (50.3\,\%) \\
    \quad 9/10 (1 mismatch) & 65 (34.8\,\%) \\
    \quad 8/10 (2 mismatches) & 23 (12.3\,\%) \\
    \quad 7/10 (3 mismatches) & 5 (2.7\,\%) \\
    \addlinespace
    CD34$^+$ dose ($\times 10^6$/kg) & 9.7 (5.4--15.4) \\
    High-risk disease & 69 (36.9\,\%) \\
    Malignant disease & 155 (82.9\,\%) \\
    \addlinespace
    \multicolumn{2}{l}{\textit{Outcomes (competing-risk encoding)}} \\
    \quad Acute GvHD grade III--IV (alive, no relapse) & 16 (8.6\,\%) \\
    \quad Relapse & 28 (15.0\,\%) \\
    \quad TRM (dead without relapse) & 62 (33.2\,\%) \\
    \quad Censored / alive & 81 (43.3\,\%) \\
    \quad Follow-up, median days & 676 (168--1604) \\
    \bottomrule
  \end{tabular}
\end{table}

\begin{table}[htbp]
  \caption{Integrated Brier Score (IBS) integrated over 6-month, 12-month, and
    18-month evaluation times (95\,\% bootstrap CI, 1\,000 resamples) for all
    models on the held-out test set ($n=29$).  Lower is better.
    Bold: best result per column.
    GvHD IBS is near zero for all models due to the rarity of events (8.6\,\%
    of the cohort) and is not a meaningful discriminator.
    All models use aggregate HLA mismatch features; per-allele ESM-2 embeddings
    require registry-scale data.}
  \label{tab:ibs}
  \centering
  \begin{tabular}{lccc}
    \toprule
    \textbf{Model} & \textbf{GvHD} & \textbf{Relapse} & \textbf{TRM} \\
    \midrule
    Cox (cause-specific)        & 0.000 (0.000--0.000) & 0.114 (0.011--0.245) & 0.202 (0.127--0.286) \\
    \textbf{Fine--Gray}         & 0.000 (0.000--0.000) & \textbf{0.110} (0.016--0.232) & \textbf{0.202} (0.132--0.279) \\
    Random Survival Forest      & 0.000 (0.000--0.000) & 0.129 (0.016--0.262) & 0.207 (0.129--0.293) \\
    DeepHit (tabular HLA)       & 0.000 (0.000--0.000) & 0.117 (0.006--0.254) & 0.269 (0.191--0.345) \\
    \bottomrule
  \end{tabular}
\end{table}

\begin{table}[htbp]
  \caption{CAPA model parameters and estimated computational cost (architectural
    design values).  Inference times estimated for a single NVIDIA T4 GPU with
    batch size of 1, based on component complexity.  These values characterise
    the architecture and have not been benchmarked on the UCI BMT cohort, which
    does not contain the allele-level data required to run the ESM-2 pipeline.
    The cross-attention network parameter count ($\sim$27.6\,M) is dominated
    by the four multi-head attention blocks operating at embedding dimensionality
    $d = 1280$ (each block: $4d^2 \approx 6.6\,\text{M}$); the post-projection
    MLP ($2560 \times 512 + 512 \times 128$) contributes $\sim$1.4\,M.
    The DeepHit survival head ($\sim$185\,K) corresponds to the shared MLP
    ($160 \to 256 \to 256$) plus three event-specific linear heads
    ($3 \times 256 \to 100$).}
  \label{tab:compute}
  \centering
  \begin{tabular}{lrr}
    \toprule
    \textbf{Component} & \textbf{Parameters} & \textbf{Inference latency} \\
    \midrule
    ESM-2 650M (frozen) & 650\,M & $<$1\,ms (cached) / $\sim$120\,ms (live) \\
    Cross-attention network & $\sim$27.6\,M & $<$5\,ms \\
    Clinical encoder & $\sim$18\,K & $<$1\,ms \\
    DeepHit survival head & $\sim$185\,K & $<$1\,ms \\
    \midrule
    Total learnable & $\sim$27.8\,M & $\sim$7\,ms (with cache) \\
    \bottomrule
  \end{tabular}
\end{table}

\begin{table}[htbp]
  \caption{Reproducibility metadata for all experiments reported in the main
    text.  All results in Table~1 and Supplementary Table~S2 can be exactly
    reproduced by running \texttt{uv run python scripts/run\_real\_baselines.py}
    from the repository root with the package versions listed here.}
  \label{tab:repro}
  \centering
  \begin{tabular}{ll}
    \toprule
    \textbf{Item} & \textbf{Value} \\
    \midrule
    Random seed (all models)   & 42 \\
    Train / val / test split   & 70\,\% / 15\,\% / 15\,\% (stratified by event type) \\
    \addlinespace
    Python                     & 3.11 \\
    PyTorch                    & 2.x \\
    lifelines                  & 0.30 \\
    scikit-survival            & 0.27.0 \\
    scikit-learn               & $\geq$1.3 \\
    numpy                      & $\geq$1.24 \\
    \addlinespace
    Bootstrap iterations       & 1\,000 \\
    CV configuration           & 5 folds $\times$ 5 repeats, stratified \\
    Baseline script            & \texttt{scripts/run\_real\_baselines.py} \\
    \bottomrule
  \end{tabular}
\end{table}

%% ── References ──────────────────────────────────────────────────────────────
\bibliographystyle{abbrvnat}
\bibliography{references}

\end{document}
